\begin{table}[htbp]
\centering
\caption{External validation of Frozen B5/P3 SGAF on two BEIR datasets not used in transfer development. TREC-COVID (50 queries) and ArguAna (1406 queries) test generalization to biomedical IR and argument retrieval domains respectively. All SGAF variants use the SciFact-fitted batch-shift detector frozen and applied zero-shot. Cost is ``negligible'' --- B5 routing and P3 smoothing operate on already-computed top-100 lists with CPU-only post-retrieval overhead ($\sim$0.01\,ms and $\sim$0.001\,ms respectively).}
\label{tab:external-validation}
\small
\resizebox{\textwidth}{!}{%
\begin{tabular}{lrrrrrrl}
\toprule
Dataset & Queries & BGE-base & Frozen B5 & Frozen P3 & $\Delta$(P3$-$B5) & $p$ & Sig. \\
\midrule
TREC-COVID & 50   & 0.6835 & 0.6835 & 0.6908 & +0.0073 & 0.47 & n.s. \\
ArguAna    & 1406 & 0.4530 & 0.4301 & 0.4312 & +0.0011 & 0.02 & * \\
\bottomrule
\end{tabular}
}

\vspace{0.5em}
\footnotesize
\textbf{Interpretation:} On TREC-COVID, B5 matches BGE-base (all 50 queries classified as generalist-fallback). P3 adds a small +0.0073 gain via rank-window smoothing but the 50-query sample lacks statistical power ($p=0.47$). On ArguAna, B5 underperforms BGE-base by $-$0.0229 --- the SciFact-trained batch-shift detector falsely activates the BGE-small specialist on some argument-retrieval batches where BGE-base is the correct choice. P3 partially compensates (+0.0011, $p=0.02$) by applying specialist-informed smoothing inside those misrouted batches, but cannot fully close the gap. This reveals a fundamental limitation: the SciFact-fitted source-shift detector does not generalize to argument retrieval.
\\
\textit{Data provenance:} \texttt{external\_validation/trec\_covid\_beir/external\_validation\_summary.csv} and \texttt{external\_validation/arguana/external\_validation\_summary.csv}. Bootstrap $p$-values (10\,000 resamples, paired).
\end{table}
